\begin{workedexample}[Worked Example: Band and wavelength]
The IEEE letter bands split the microwave spectrum: L (1--2 GHz), S (2--4),
C (4--8), X (8--12), Ku (12--18), K (18--27), Ka (27--40). A half-wave dipole
scales as $\ell = \lambda/2 = c/(2f)$: at S-band $2.45\ \text{GHz}$,
$\ell = (3\times10^8)/(2\cdot2.45\times10^9) = 61\ \text{mm}$; at Ka-band
$30\ \text{GHz}$ it is only $5\ \text{mm}$.
\end{workedexample}

\begin{keytakeaways}
\begin{itemize}[leftmargin=1.4em,itemsep=2pt]
\item IEEE bands (L, S, C, X, Ku, K, Ka) label microwave ranges.
\item Antenna size scales with $\lambda = c/f$; higher bands mean smaller antennas.
\item Band choice trades bandwidth/size against rain and atmospheric loss.
\end{itemize}
\end{keytakeaways}

\begin{exercises}
\item Which IEEE band contains $10\ \text{GHz}$?
\item Half-wave dipole length at $1\ \text{GHz}$?
\item Do higher bands give larger or smaller antennas at the same electrical size?
\end{exercises}

\furtherreading{Pozar~\cite{pozar}; IEEE Std~521~\cite{ieee521}.}
